\section{Limitations}
\label{sec:limitations}

\paragraph{One generator family.}
Everything here is AMLworld. The tail and warm-up are reproduced on LI-Small; on HI-Medium the
warm-up, the budget capture, the share of unobservable rings, the threshold dependence of the
partition gain and the direction (not the size) of the scope effect reproduce, from a detector
that missed its adequacy bar and two usable seeds; the transferred detector's seeds-only rings
are the only HI-Medium method numbers we report under transfer. All three releases come from one simulator;
nothing establishes how the findings carry to other generators or to real transaction data. Method
comparisons rest on 174 truth rings and 5 detector seeds, and only one of five pre-registered
comparisons survives correction at that power.
% source: docs/FINDINGS.md (A6); docs/paper/generalisation.md §3

\paragraph{External baselines.}
We score two external group-level methods (CSGM over five random two-institution splits; TAML,
three seeds, converted to rings) and two families of community detection under the protocol.
TAML is scored with protocol differences we could not remove: its features aggregate all ten
days, and its graphs omit about 42\,\% of in-window ring transactions. Leiden was planned but not
run, and \citet{chai2023} was not run.
% source: results/tables/v40_external_baselines.json (runs, notes); results/notes/p3_baselines.md

\paragraph{The edge-versus-ring claim is conditional.}
It is measured at one seed budget, across three graph neural detectors from edge F1 0.596 to
0.715, on one dataset. Below that range it does not hold (GFP-XGBoost, HI-Medium), and detectors
above it were not tested.

\paragraph{A ring-level leak in the expansion policy's training data.}
The policy's cross-scored inputs hold out a day per fold, but rings span days, so a fold model
is supervised on some edges of a ring and then scores that ring's held-out-day edges. Between
65.0\,\% and 92.4\,\% of each fold's held-out rings are affected. The measured effect is small
--- the gap in ROC-AUC between affected and unaffected ring edges ranges from $-0.008$ to
$+0.036$ across folds, and unaffected rings are still scored at AUC 0.958--0.996 --- and the
leak touches training inputs only: test scores come from a detector trained on the train
split alone. Removing the affected feature group changes primary-window ring-F1 by 0.008,
which bounds what a leak-free version could change on HI-Small; on HI-Medium the same ablation
moves validation ring-F1 by 0.026 in the opposite direction, so the bound does not transfer
between releases.
% source: docs/FINDINGS.md (B17); results/tables/v33_ring_fold_leakage.csv, v33_leak_inflation.csv, v33_policy_ablations.csv
We re-ran the policy with its inputs folded by ring, so that no fold model sees any edge of a
ring it scores. The registered check on the mean ring-F1 difference passed ($+0.0028$, inside
$\pm 0.008$), but the paired interval, $[-0.0137, +0.0211]$ ($p=0.77$), is too wide to call the
two equivalent, so on ring-F1 the comparison is inconclusive. The leak did inflate precision:
with leak-free inputs, ring precision falls from 0.838 to 0.784 ($-0.055$ [$-0.089$, $-0.027$])
and added-member precision from 0.926 to 0.819, while ring recall rises by 0.022. The full
pipeline behaves the same way: ring-F1 0.589 with either input ($-0.000$ [$-0.017$, $+0.019$]),
added-member precision 0.926 against 0.817. Where this paper reports added-member or ring
precision, the leak-free value is given first. These re-runs were added after the main results
and are exploratory.
% source: docs/deep-analysis.md §8 (verdict) and §15; results/tables/v39_ringfold_report.csv (paired: ring_f1, ring_precision, ring_recall, added_member_precision); v45_ringfold_full_precision_paired.csv (full_ringfold_vs_full); v45_posthoc_ring_ci.csv; paper/PREREG_AMENDMENT_2026-11-14.md

\paragraph{Analyses added after the main results.}
The precision-aware paired test, the window-matched budget, the component-level scope and the
reallocation analysis were run on a test window that had already been examined many times.
Their designs and bars were written down before each was run, but they are sensitivity
analyses, not confirmatory tests. None changes a reported point estimate.
% source: docs/FINDINGS.md (section G preamble); docs/deep-analysis.md (header and caveats)

\paragraph{Published-number comparison.}
The only published figure we compare against is MEGA-GNN's edge F1
(Section~\ref{sec:results-swap}).

\paragraph{Evaluation defects we found.}
Each changed a number we had reported, and each is corrected here: bootstrap resamples that
collapsed repeated rings (intervals about a quarter too narrow); matcher ties broken by file
order; detector scores that depend on the scoring batch through batch-level causal sampling, and
the look-ahead that sampling allows (Section~\ref{sec:data}); a ring-level leak through
day-folded policy inputs (above); a global quantile that hands the seed budget to one day; and
saturated scores that leave a strict quantile cut empty (Section~\ref{sec:findings}).
% source: results/notes/interval_consistency.md; matcher_tiebreak.md; p7_genagg.md (scoring batch, look-ahead); FINDINGS B17, C11; findings saturation
